% Pursuing Stacks (cotes 134-2 to 134-8): archive cote 134-2, site page 4 (PDF page 5).
% provisional contents. Draft compared: none.
% AI-transcribed by Claude Opus 5.5 (claude-opus-5-5, adaptive thinking, effort high), Message Batches API, batch ps1 (msgbatch_01Rv5hCaNTtUt4JHx1yqoHPL), 2026-11-12; reviewed against both supplied scan halves; see CHECK_DECISIONS.tsv.
\scanpage{134-2/4}
\over{\unsure{PURSU}}{PURSU}ING STACKS\par

( A la \over{p}{P}oursuite des Champs )\par

First episode : THE MODELIZING STORY\par
% NOTE: a hand-drawn curved line runs from under "episode" around the opening parenthesis of the next line
(Histoire de Modèles )\par

\uline{C\over{\unsure{o}}{O}ntents} \hand{,}\par

% NOTE: a hand-drawn line runs from the left margin under "I The Take-off"
\hand{I} \uline{The \over{t}{T}ake-off} (a letter to Daniel Quillen)\par
1. The importance of innocence.\par
2. A \deleted{quick}\inserted{short} look into purgatory\par
3. \hand{``}Fundamen\unsure{tt}al \hand{$\infty$}- groupoids\hand{''}as objects of a "model category" ?\par
% NOTE: a hand-drawn line runs from the left margin between items 3 and 4
4. A bit of ordering in the mess of \over{"}{``}higher order structures\hand{''}.\par
\hand{\textcircled{5.}} Jumping over the ab\deleted{\unsure{y}}yss !\par
6. T\over{\unsure{n}}{h}e topological model: hemispheres building up the (tentative)\par
\hand{``}universal \hand{$\infty$}- (co) groupoïd\hand{''}.\par
7. Gluing hemispheres: the "standard\over{"}{''} amalgamations.\par
% NOTE: a hand-drawn line runs from the left margin between items 7 and 8
8. Description of the universal \uline{prim\over{\unsure{ary}}{itive}}\hand{/}structure.\par
\hand{\textcircled{9.}} The main inductive step: just add coherence arrows ! \inserted{And the} \deleted{The} \xout{\unsure{storyxxofx}} \xout{\unsure{THEXIRRELEVANT}}\par
\inserted[typed]{abridged story} \xout{ambiguity} of an \hand{(}unescapable and irrelevant\hand{)} ambiguity ...\par
10. Cutting down redundancies - or\hand{:}"l'embarras du choix".\par
11. Returning to the topological model\deleted{\unsure{: getting}} \hand{(}\inserted{\unsure{i.e.|e.g.}} \over{t}{T}he \inserted[typed]{canonical functor} \xout{\unsure{geometricxrealization}}\par
\xout{\unsure{for(oo-groupoids))x}} fr\over{\unsure{o}}{om} spaces to "\unsure{$\infty$}-groupoids"\hand{)}.\par
12. About replacing spaces by objects of a "model category" \hand{.}\par
\hand{\textcircled{13.}} \inserted[typed]{An urgent} \xout{Reflection} reflection on proper names: "\over{s}{S}tacks" and "coherators".\par
% NOTE: a hand-drawn double line runs across the whole page here
\deleted{\illegible{1}}\par
II \hand{\uline{Test categories} \uline{and} \uline{test functors}.}\par
14. \deleted{Homotopical algebra versus "topology"} The \inserted[typed]{unnoticed \deleted{breakdown}} \inserted{failure} \xout{\unsure{essentialxinadequacy}}\par
(of the present foundations of topology for expressing topological\par
intuition).\par
15. Overall review \hand{(}on\hand{)} \inserted[typed]{standard} descriptions of homotopy types. \deleted{by models.}\par
16. \deleted{The main motivation: stacks over topoi.} Stacks \inserted{over topoi} as the unifying\par
concept for homotopical and cohomological algebra.\par
1\over{7}{5}.\hand{bis} Categori\over{\unsure{a}}{e}s as models for homotopy types. F\over{\unsure{u}}{i}rst glimpse upon an\par
"impressive bunch" \hand{(of modelizers)}.\par

\mspage{\hand{4}}
